\subsection{Project Background and Motivation}

Conventional communication systems are firmly rooted in Shannon's mathematical theory of communication, where the main design objective is accurate bit delivery through noisy channels \cite{shannon1948, shannon1948part2}. This objective has driven decades of successful progress in source coding, channel coding, modulation, and multiple-access schemes. However, an increasing number of emerging intelligent applications, such as human--machine dialog systems, distributed artificial-intelligence (AI) collaboration, autonomous driving control loops, and edge semantic sensing, do not strictly require exact bitwise reconstruction if the underlying meaning of the message remains useful for the downstream task. This conceptual gap between syntactic fidelity and semantic usefulness has motivated the recent resurgence of semantic communication, in which the design objective is redefined from bit-level fidelity to meaning-level fidelity \cite{carnap1954, bao2011semantic}.

Recent research has shown that deep neural networks can learn task-oriented semantic representations and perform end-to-end semantic transmission over realistic channel models \cite{xie2021deepsc, lu2024survey}. In particular, Transformer-based semantic transceivers have become a mainstream design choice thanks to their strong contextual modelling capability, their parallelizable structure, and their compatibility with pretrained language priors \cite{vaswani2017attention}. Despite the promising progress reported in the literature, several deployment-oriented issues remain unsolved. Low-SNR and fading conditions still cause severe semantic degradation; model components that appear theoretically beneficial may become unstable in practice when interacting with channel noise and other components; external post-processing models driven by language priors may introduce additional error rather than effective correction; and the behavior observed under an AWGN benchmark does not necessarily transfer to Rayleigh fading scenarios that are closer to real wireless propagation.

This project is motivated by a very practical engineering question: can one construct a semantic transceiver that is not only theoretically improved, but also experimentally validated through a rigorous iterative process whose outcomes, whether positive or negative, are reproducible and traceable? The emphasis here is not merely on ``adding more neural modules'' on top of an existing pipeline. Instead, this project aims to establish a full hypothesis--diagnosis--redesign--validation loop that identifies both effective design choices and ineffective ones.

From an undergraduate research perspective, this motivation has two very practical implications. First, the project should remain reproducible and implementation-feasible on a single GPU workstation, so that each design decision can be traced back to concrete code changes and to measurable outcomes in saved evaluation files. Second, the final report should not merely highlight improvements but also explicitly document negative findings and their root causes, because engineering decisions in real systems are typically constrained by robustness, interpretability, latency, and memory footprint rather than by isolated best-case numbers.

In this context, the project combines model-side architectural design, training-side objective tuning, and inference-side strategy optimization. This integrated view is necessary because semantic communication quality is jointly affected by representation learning behavior, channel perturbation statistics, and decoding policy interactions. A single-module improvement may appear promising in isolation yet fail after being integrated with other modules or evaluated under a different channel regime. In this project, each design decision is therefore treated as a testable hypothesis rather than as an assumed improvement.

Recent semantic communication studies also suggest that the motivation for this project should not be understood only as a replacement of conventional bit-level communication. A more accurate interpretation is that semantic communication complements conventional communication by introducing a new optimization layer above physical-layer reliability. In a conventional system, successful delivery is normally judged through bit error rate, symbol error rate, packet error rate, or throughput. These measures are still essential when the receiver must reproduce the exact transmitted message. However, in many intelligent communication scenarios, the receiver is not a passive data sink but a task-performing agent. For example, if a downstream classifier, dialogue agent, or control module can still make the correct decision from a compressed or partially corrupted representation, then exact bit-level reconstruction may be unnecessarily strict. This motivates a task-aware and meaning-aware design philosophy in which communication resources are allocated according to semantic importance rather than uniform syntactic fidelity.

This shift is especially relevant to future wireless systems because the expected traffic pattern is becoming increasingly machine-centric. Edge devices, autonomous agents, sensors, and embedded AI systems may exchange not only raw data but also intermediate features, intents, instructions, and task-specific semantic states. Under bandwidth limitation, latency constraint, and energy constraint, transmitting compact semantic representations can be more efficient than transmitting all raw symbols and reconstructing them perfectly. Therefore, a semantic transceiver should be evaluated not only by whether it can reproduce the sentence exactly, but also by whether it preserves enough meaning for the receiver to perform the intended reconstruction or task. This is consistent with the broader post-Shannon research trend, where semantic extraction, semantic encoding, adaptability, and reliability are treated as central challenges rather than secondary details \cite{liu2024survey_dcn, zhang2024semantic_nso}.

For this project, the most important motivation is robustness rather than compression alone. Text transmission through a neural semantic transceiver is vulnerable to channel noise in a different way from traditional coded communication. In a traditional coded system, errors are often localized at the bit or symbol level and can be analyzed through classical coding theory. In a neural semantic transceiver, however, channel perturbation affects continuous latent vectors. A small perturbation in the latent space may produce either a minor lexical variation or a large semantic shift after autoregressive decoding. This nonlinear error propagation makes low-SNR behavior especially important. The experimental focus on \SI{-10}{dB} and \SI{-5}{dB} therefore reflects the practical operating region in which robust semantic modelling is most strongly tested, not simply an arbitrary extension of the SNR grid.

Another motivation is the need to evaluate apparently reasonable neural-network improvements with caution. Semantic communication papers often report performance gains by adding attention modules, semantic losses, or pretrained language-model components. However, these modules may interact in unexpected ways. A semantic loss may encourage global sentence similarity but reduce token-level reconstruction. A language model may improve fluency but alter the original meaning. A channel-aware encoder may help representation learning but still provide insufficient information for decoder-side recovery. Therefore, each added module in this project is tested against a controlled baseline and judged by its measured effect under the same evaluation protocol.

\subsection{Research Questions and Objectives}

To make the investigation testable and reproducible, the project is organized around four explicit research questions that drive the design and experimental program:

\begin{enumerate}
    \item Can channel-quality-aware attention mechanisms, placed inside both the encoder and the decoder, improve semantic reconstruction quality under low SNR conditions?
    \item Does introducing a semantic-regularization term into the loss function consistently improve final communication quality, or does it instead introduce gradient conflict with token-level cross-entropy?
    \item Can a pretrained language model be used for semantic post-refinement in a way that provides reliable net gains on top of an already well-trained decoder?
    \item If learnt post-refinement proves unstable, can a training-free decoding strategy, such as beam search, offer a better engineering trade-off in terms of accuracy, cost, and determinism?
\end{enumerate}

Guided by these four questions, the concrete engineering objectives of the project are:

\begin{itemize}
    \item Build and verify a faithful baseline semantic communication transceiver for text transmission, compatible with the literature and reproducible on commodity hardware.
    \item Implement three distinct innovation modules covering model design, training objective, and inference refinement, each controlled by a clean command-line flag.
    \item Evaluate all variants under both AWGN and Rayleigh channels on an identical SNR grid and with identical metric definitions, so that every comparison is fair and well-scoped.
    \item Quantify effectiveness boundaries and failure mechanisms, not only best-case metric values, and report both positive and negative outcomes with clear attribution.
    \item Translate the observed evidence into concrete, deployment-oriented design recommendations that can guide follow-up final-year projects and early-stage research work.
\end{itemize}

To ensure that conclusions are objective, all model variants are evaluated under exactly the same data split, SNR grid, decoding termination criteria, and metric implementations. The project therefore explicitly emphasizes methodological fairness: when one module is added or removed, all other conditions remain fixed, and the resulting evaluation files are committed without manual cherry-picking. This principle is particularly important in semantic communication studies because several commonly used metrics, including BLEU, BERTScore, and sentence similarity, have different sensitivities to local and global errors, and inconsistent protocols could easily produce misleading conclusions in favor of a preferred variant.

\subsection{Contribution Summary}

The main contributions of this project can be summarized as follows:

\begin{itemize}
    \item A reproducible full-stack semantic communication pipeline is completed, including data preprocessing, model training, checkpointing, evaluation, visualization, and comparison scripts, with every step documented as command-line instructions in the appendix.
    \item The SNR-adaptive attention idea is extended from encoder-only usage, which is the dominant form reported in the literature, to decoder cross-attention, enabling explicit channel-conditioned memory retrieval at the reconstruction stage.
    \item A dynamic semantic-loss weighting mechanism is introduced to replace fixed weighting, so that the semantic-consistency term is de-emphasized under poor channel conditions and gradually re-emphasized under cleaner channels.
    \item A confidence-driven BERT masked-refinement strategy is implemented to replace earlier unconstrained GPT-style corrections, focusing edits on low-confidence positions and restricting global rewriting behavior.
    \item A beam search decoding strategy with beam width five is integrated and evaluated as an inference-stage method that requires no retraining.
    \item Negative experimental results are reported and analysed, including the weak performance of the semantic-loss component and the instability of LLM-based refinement.
\end{itemize}

In addition to the implemented modules, the project emphasizes controlled comparison across different design choices. Each major modification is evaluated under the same data split, SNR grid, channel setting, and metric definitions where applicable. This allows the effects of SNR-aware attention, semantic-loss weighting, LLM refinement, and beam search to be compared more directly. The resulting discussion therefore includes both performance improvements and negative results, rather than reporting only the best-performing configuration.

\subsection{Structure of This Report}

The remainder of this report is structured as follows. Section 2 presents the technical body, including literature context, the baseline architecture, the initial innovation hypotheses, the diagnosed limitations of those hypotheses, the resulting iterative redesign process, and the quantitative results of each configuration on AWGN and Rayleigh channels. Section 3 focuses on analysis and discussion, where the mechanisms behind the observed gains and failures are interpreted and connected with standard wireless-channel behavior and deep-learning optimization theory. Section 4 summarizes the conclusions and proposes concrete future work directions along both scientific and engineering axes. The references section then lists all external works cited in the main text, ordered by citation order.

The appendices provide reproducibility resources, including the complete command pipelines used for training and evaluation, and the mapping between experimental variants and saved evaluation files under the \texttt{checkpoints/} tree. Together, the main text and the appendices are designed to let a future student continue this line of work without reconstructing the entire experimental setup from scratch, and to allow examiners to verify any numerical claim by running the listed commands directly.
